\chapter{基于中心化的逻辑时间戳分配节点的混合隔离级别的分布式事务协议的设计}
本章将基于上一章中所提到的建模、公理和隔离级别公约，基于中心化的逻辑时间戳分配节点，实现一个分布式复制数据库的原型。
在实现的过程中，我们设计的协议需要满足上一章节中所提到的隔离级别公约，向客户端提供一致性和隔离性的保证的同时，给出相应的具体操作的接口。

在设计协议时，我们需要重点关注三个操作：事务在开始时应当选择哪种快照；如何决定事务之间的逻辑上的先后顺序；如何处理事务之间的冲突。

\section{时间戳和时间戳生成}
先发生(happened-before)是由Lamport\cite{happened-before}提出的分布式系统中一个重要的概念，用于描述事务之间的逻辑顺序。
其目的是解决分布式场景下 “如何判断两个事件的先后顺序” 这一基础问题。
在事务内部，由于每个事务的请求独占某一个节点(replica)，事务内部操作的逻辑顺序是天然存在的。
事务之间，由于事务请求分散在不同的节点执行，事务之间的逻辑顺序不直接存在，需要借助其他手段来决定。

通过时间戳预言机(Timestamp Oracle，TSO)\cite{tso1}来解决确定事务间逻辑上的先后顺序问题。
TSO是分布式系统中一种中心化的服务，是为事务或事件分配全局唯一、单调递增的时间戳的核心服务或者组件。
它的核心作用是解决分布式环境下缺乏统一时钟的问题，为事务提供统一的时间基准，确保系统能实现强一致性（如线性一致性或外部一致性）。
在众多分布式数据库中也有使用时间戳预言机的，
像 Google Spanner 使用 TrueTime API 作为 TSO\cite{spanner}，
TiDB 则依靠 PD（Placement Driver）服务充当 TSO 角色\cite{TiDB}。
它们生成的时间戳用于确定分布式事务的执行顺序，实现一致性读写，让不同节点读取到的数据保持统一视图。

TSO 用于决定 AR 上的全序关系，每个事务在提交的时候必然持有一个全局唯一的自增的提交时间戳(commit timestamp, cts)，
若\Txc[1] 在\Txc[2] 之前提交，那么 $T_1.\text{cts} < T_2.\text{cts}$。
VIS不是一个全序关系，但是它是AR的子集，TSO在决定VIS上的偏序关系提供帮助。
时间戳也是事务开始的标志，每个完整完成的事务在持有提交时间戳cts的同时还有开始时间戳(start timestamp, sts)。
如果一个$T_1 \xrightarrow{\text{VIS}} T_2$ 意味着 $ T_1.\cts \leq T_2.\sts$。
接下来的章节中，我们将直接使用时间戳的比较关系来反映事务之间的逻辑顺序和它们直接的(AR,VIS)关系。

对于事务\Txc[1]和\Txc[2]，如果$T_1.sts < T_2.sts  \wedge T_1.cts > T_2.cts$，
也就是事务\Txc[1]在事务\Txc[2]之前开始，但\Txc[1]在\Txc[2]之后提交，那么这两个事务是并行的(concurrent)，
记作$T_1 || T_2$。

算法~\ref{alg:timestamp-oracle} 展示了一个简化的时间戳预言机的核心操作。
TSO 维护一个单调递增的计数器 \textbf{counter}，
通过两个基本操作实现时间戳的分配：
\textsc{Tick()} 操作用于读取当前的时间戳值，
而 \textsc{Tock()} 操作则将计数器递增并返回新的时间戳。
这种设计确保了系统中所有时间戳的全局唯一性和单调性，
为分布式事务提供了统一的逻辑时钟基础。


\begin{algorithm}[H]
\caption{Timestamp Oracle}
\label{alg:timestamp-oracle}
\begin{algorithmic}[1]
\State \textbf{counter}: monotonically increasing timestamp counter
\Procedure{Tick()}{}
    \State \Return $\text{counter}$ \Comment{return current timestamp}
\EndProcedure
\Procedure{Tock()}{}
    \State $\text{counter} \gets \text{counter} + 1$
    \State \Return $\text{counter}$ \Comment{increment and return new timestamp}
\EndProcedure
\end{algorithmic}
\end{algorithm}

\section{操作语义的设计}
本节，我们将从RA隔离级别出发，向SER一步一步推导，设计各个隔离级别的操作语义，并证明其正确性。

\subsection{RA(Read Atomic)}

实现RA隔离级别需要保证内部一致性和外部一致性公理(INT和EXT)。
在事务执行时，需要选择一个快照(snapshot)，这个快照由多个事务的共同作用决定。
在事务的执行过程中，写操作的操作结果需要预先写入一个缓存；
读操作从缓存中读取数据，如果数据不存在于缓存中，则从快照中读取数据。
在提交时，写入操作需要保证要么全部成功提交，要么全部不提交，即要么所有写操作都成功，要么所有写操作都不成功。

\begin{equation}
\begin{aligned}
S &= \langle T_1,\dots,T_n\rangle
\quad\text{s.t.}\quad
\forall\, 1 \le i < j \le n.\ T_i\  \text{AR} \  T_j \\
C_1(k,v) &\triangleq (v = T.\text{buffer}(k))\\
C_2(k,v) &\triangleq S(k,v)\\
C_3(k,v) &\triangleq (v = \bot)\\
Read(k,v) &\Longleftarrow C_1(k,v) \lor C_2(k,v) \lor C_3(k,v)\\
          &\land \neg(C_1(k,v) \land C_2(k,v))\\
          &\land \neg(C_1(k,v) \land C_3(k,v))\\
          &\land \neg(C_2(k,v) \land C_3(k,v))
\end{aligned}
\end{equation}

在实现时，每个节点会持有变量 deps，这是一个局部变量，用来记录该节点之前接收到哪些事务的提交时间戳。
和MVCC相似，它由两部分组成，minDep和depSet，minDep是一个时间戳，表示小于等于该时间戳的事务都已被该节点接收，
depSet是时间戳大于minDep的已经被该节点接收的事务的时间戳的集合。deps意味着哪些事务是当前节点可见的。
\begin{equation}
    \forall T_1, T_2.\ \ T_2 \in T_1.\text{deps} \implies T_2 \xrightarrow{\text{VIS}} T_1
\end{equation}

对于RA级别，对开始的时间戳没有特别的限制，事务可以选择deps中最大的时间戳作为自己的开始时间戳，即$T_i.\text{sts} = \max(\text{deps})$。
在提交时，节点向TSO发送一个\tock 请求，申请一个新的时间戳作为自己的提交时间戳。
RA隔离级别下，只需要关注内部一致性和外部一致性，无需考虑事务之间的冲突，时间戳更大，也就是更晚提交的事务会覆盖更早提交的事务。
为了保证和其他隔离级别语义的一致性，在提交时间，节点也需要将当前deps一并提交。

RA隔离级别较低，只需要保证事务的提交或者回滚操作是原子的，即要么全部提交，要么全部回滚，就可以保证该事务的正确性。

RA 隔离级别仅保证事务更新的原子可见性，
即其他事务要么看到该事务的全部写入，要么完全看不到。
由于 RA 不对事务之间的可见性关系施加传递性、前缀一致性或冲突检测等额外约束，
因此几乎所有并发异常都可能发生，
包括因果违背（Causality Violation）、丢失更新（Lost Update）、长分叉（Long Fork）和写偏斜（Write Skew）。
这使得 RA 成为最弱的一致性级别，但同时也具有最低的性能开销和最高的并发度。

\subsection{CC(Causal Consistency)}
CC 相较于RA，增加了传递可见性公理(TransVis)，如果事务\Txc[1]对\Txc[2]可见，
那么对\Txc[1]可见的任何事务对\Txc[2]也必须可见，即强调 VIS关系的传递性。

在实现时，每个新任务的deps有如下性质：
\begin{equation}
    T.\text{deps} = \bigcup\nolimits_{S \xrightarrow{\text{VIS}} T} (S.\text{deps} \cup \{S\})
\end{equation}
CC隔离级别下，节点开启事务的前提条件是deps中的所有事务的deps都是deps的子集，即所有对deps中事务可见的事务对当前快照必须可见：
\begin{equation}
    \forall S \in T.\text{deps}.\ S.\text{deps} \subseteq T.\text{deps}
\end{equation}

要证明CC隔离级别的正确性，需要通过deps实现VIS关系的传递性。我们使用反证法进行证明。
假设存在\Txc[1]、\Txc[2]、\Txc[3]，满足$T_1 \xrightarrow{} T_2$, $T_2 \xrightarrow{} T_3, T_1 \nrightarrow T_3$，
即$T_1 \in T_2.\text{deps}, T_1 \in T_2.\text{deps}$，从而$T_1 \in T_3.\text{deps}$，得到 $T_1 \xrightarrow{} T_3$，
与假设相违，证得deps实现了VIS关系的传递性。

可以补充到：对于任意的事务\Txc，VIS关系在$T.\text{deps}$下的限制是自反传递闭包。
\begin{equation}
    \text{VIS} \upharpoonright (T.\text{deps}) = \left( \text{VIS} \upharpoonright (T.\text{deps}) \right)^*.
\end{equation}

在CC级别下，事务开始时是需要检查deps中的所有事务的deps是否都是当前deps的子集，这种状态可以被定义为因果交付(Causal Delivery)。
只有当前快照满足因果交付的条件，才能够执行CC隔离级别的事务。

CC 隔离级别通过强制 VIS 关系的传递性，
有效消除了因果违背异常，
保证了事务之间的因果依赖关系能够被正确传播。
然而，由于 CC 并不禁止对同一对象的并发写操作，
丢失更新（Lost Update）异常仍然被允许。
同时，CC 也不要求事务看到一致的历史前缀，
因此长分叉（Long Fork）和写偏斜（Write Skew）异常同样可能出现。

\subsection{PC(Prefix Consistency)}
PC隔离级别在CC的基础上，增加了前缀一致性公理(Prefix Consistency)，即：
\begin{equation}
 \forall T', S \in \mathcal{T}.\ T' \xrightarrow{\text{AR}} S \land S \xrightarrow{\text{VIS}} T \Rightarrow T' \xrightarrow{\text{VIS}} T
\end{equation}
前缀一致性公理实际上是对传递可见性公理的加强，它强调了如果一个事务对另外一个事务可见，那么它所覆盖的事务对另一个事务也必须可见。
我们在前文中提到，TSO生成的时间戳是单调递增的，提交时间戳的大小关系反映事务在AR下的全序关系；
而deps集合意味着当前事务可见的事务的集合，对于其中具有最大提交时间戳的事务\Txc[max]，有
\begin{equation}
    \forall T'\in \mathcal{T}.\quad T'.\text{cts} \le T_max.\text{cts} \implies T' \xrightarrow{\text{AR}} T_{max}
\end{equation}
又因为\Txc[max] 对当前事务可见，所有小于等于\Txc[max]提交时间戳的事务都应当被当前事务看到。
从而要实现前缀一致性，需要满足：
\begin{equation}
    \forall T \in \mathcal{T}, T.\text{cts} \leq t.\  T.\text{cts} \in \text{deps}
\end{equation}
对于一个指定的时间戳ts，如果快照$\mathcal{S}$ 满足：
% 所有小于等于ts的事务都被接收来构成这个快照
\begin{equation}
\begin{aligned}
    S_\text{ts} &= \langle T_0,T_1,\dots,T_\text{ts}\rangle \\
      &\quad\text{s.t.}\quad \forall 0 \leq u \le \text{ts}, T_\text{i}.\text{cts} + 1 = T_\text{i+1}.\text{cts}
\end{aligned}
\end{equation}
即，所有小于等于ts的事务都被接收来构成这个快照，那么这是快照就满足完全交付(Total Delivery)。完全交付对前缀一致性公理的实现无需额外证明。

开始事务时，和RA隔离级别一样，可以选择deps中最大的时间戳作为自己的开始时间戳，或者向客户端发送\tick 请求，获取一个更新的时间戳作为开始时间戳。
相对于RA隔离级别，PC隔离级别在事务开始时需要保证所有提交时间戳小于等于开始时间戳的事务都被接收，满足完全交付。

PC 隔离级别通过完全交付约束，
确保每个事务看到的历史是全局 AR 序列的一段连续前缀，
从而有效防止了长分叉（Long Fork）异常。
前缀一致性公理蕴含了传递可见性，
因此 PC 同样能够消除因果违背异常。
然而，由于 PC 不禁止在同一快照上对不同对象进行并发写入，
写偏斜（Write Skew）异常在 PC 中仍可能发生。
此外，PC 也不检测对同一对象的并发写冲突，
丢失更新（Lost Update）异常同样被允许。

\subsection{PSI(Parallel Snapshot Isolation)}
PSI 隔离级别相较于CC隔离级别增加了无冲突公理，无冲突公理确保如果两个事务存在写冲突，则它们
之间必须存在可见性关系，不能彼此并发而互不可见，防止更新丢失。
这是我们第一次在隔离级别的操作语义设计中，涉及到冲突解决。

PSI隔离级别同样建立在CC隔离级别之上，和PC隔离级别不同的是，PSI的加强发生在提交阶段。
在操作语义的设计上，我们关注事务$T$在其开始时间戳$T.\text{sts}$与提交时间戳$T.\text{cts}$之间的时间区间内，
由于这段时间的事务提交对当前事务是不可见的，不应出现与当前事务\Txc 发生写写冲突且对\Txc 不可见的提交事务。
WSet(\Txc )为事务$T$的写集合(WriteSet)，PSI的核心约束可以形式化为：
\begin{equation}
    \forall T' \in \mathcal{T}.\ T'.\text{cts} \in (T.\text{sts}, T.\text{cts}) \Rightarrow \text{WSet}(T') \cap \text{WSet}(T) = \varnothing
\end{equation}
即，在当前事务$T$的可见快照之后、提交时间戳之前，不允许出现对相同键进行写入并成功提交的其他事务，从而实现 First-Committer-Wins 风格的写写冲突规约。

面向~\ref{eq:consistency}中的NoConflict公理，上述约束实现了它的逆否命题，从而保证了无冲突公理的正确性。
\begin{equation}
\mathsf{NoConflict}(T) \equiv \forall T' \in \mathcal{T}.
\ T' \centernot{\xrightarrow{\text{VIS}}} T \Rightarrow \text{WSet}(T') \cap \text{WSet}(T) = \varnothing \lor T' \centernot{\xrightarrow{\text{AR}}} T
\end{equation}

PSI 隔离级别通过无冲突公理，
强制要求对同一对象进行写入的两个事务必须通过 VIS 确定先后顺序，
从而有效排除了丢失更新（Lost Update）异常。
结合从 CC 继承的传递可见性，
PSI 同样能够消除因果违背异常。
然而，PSI 不强制要求完全交付，
不同事务在不同节点上可能看到不一致的历史前缀，
因此长分叉（Long Fork）异常仍然可能发生。
同时，由于 PSI 不检查除了写冲突以外的其他冲突，允许事务在同一快照上对不同对象进行并发写入，
写偏斜（Write Skew）异常也被允许。

\subsection{SI(Snapshot Isolation)}

SI隔离级别是对PC和PSI隔离级别的析取，在CC隔离级别的基础上既要满足PC隔离级别的完全交付，又要满足PSI隔离级别的无冲突公理。
从理论模型上看，SI 同时满足 $\mathsf{Prefix}$ 和 $\mathsf{NoConflict}$ 两个公理，
能够同时防止丢失更新和长分叉两类异常，在一致性和性能之间取得了较好的平衡。

在操作语义的实现上，SI 隔离级别需要同时满足 PC 和 PSI 两个隔离级别的约束条件。
具体而言，事务开始时需要满足完全交付条件，即所有提交时间戳小于等于开始时间戳的事务都必须被节点接收；
同时事务提交时需要满足无冲突约束，即在事务的执行期间不允许有对相同键进行写入的其他事务成功提交。

SI 隔离级别通过结合完全交付和无冲突两个约束，
在实现上保证了事务的快照一致性和写冲突检测，
从而提供了比 PSI 和 PC 更强的一致性保证。
然而，由于 SI 仍允许事务在同一快照上对不同对象进行更新，
写偏斜（Write Skew）异常仍然可能出现。
要彻底消除写偏斜，需要进一步提升到 SER 隔离级别。

\subsection{SER(Serializability)}
SER隔离级别解决了写偏斜异常，让事务执行的效果与串行化执行一致。
相较于SI隔离级别，SER隔离级别通过要求在提交阶段进行相较于SI中写写冲突检查更严格的读写冲突(WR-Conflict)检查，从而进一步增强了事务的隔离性。
SER是SI的超集，通过要求在提交阶段进行更严格的读写冲突检查，实现了真正的NoConflict公理，从而确保事务执行的效果与串行化执行一致。
即：
\begin{equation}
\begin{aligned}
&\mathsf{NoConflict}(T) \equiv \forall T' \in \mathcal{T}.
\ T' \centernot{\xrightarrow{\text{VIS}}} T \Rightarrow T \bowtie T' \lor T' \centernot{\xrightarrow{\text{AR}}} T \\ 
% 补充定义bowtie
&T \bowtie T' \equiv (\text{WSet}(T) \cup \text{RSet}(T)) \cap (\text{WSet}(T') \cup \text{RSet}(T')) \neq \varnothing
\end{aligned}
\end{equation}

值得注意的是，SER 并不强制事务以串行方式执行，而是允许事务并发执行。
其核心约束在于：
两个并发事务的读写集合不能存在交集，即它们不能访问相同的数据对象。
只有当两个事务的读写集合完全不相交时，系统才允许它们并发执行。
这种冲突检测机制确保了并发执行的结果等价于某个串行执行顺序，从而实现了可串行化保证。

\section{本章小结}
本章在上一章形式化建模与隔离级别公理体系的基础上，围绕中心化时间戳预言机（TSO）设计并实现了一个分布式复制数据库的事务协议原型。
通过 TSO 提供的全局单调递增时间戳，我们统一了事务的开始时间戳与提交时间戳，构建了基于 (AR, VIS) 的逻辑时间顺序，
并使用 \textit{deps} 结构刻画节点本地可见的事务集合，为后续各类隔离级别的操作语义奠定了统一的时钟与可见性基础。

本章从 RA 隔离级别出发，逐步增强一致性约束，依次给出了 CC、PC、PSI、SI 与 SER 的操作语义：
RA 仅要求内部一致性与外部一致性，提供原子可见性但几乎允许所有并发异常；CC 在此基础上引入传递可见性，使因果依赖能够正确传播，
从而消除了因果违背异常；PC 进一步要求完全交付，保证事务观察到的是全局 AR 序列的连续前缀，
有效防止长分叉异常；PSI 通过无冲突公理在提交阶段检测写写冲突，采用 First-Committer-Wins 规则消除丢失更新异常；
SI 同时满足完全交付与无冲突两个约束，在保证快照一致性与写冲突检测的同时，综合消除了丢失更新与长分叉两类异常；
最终，SER 通过更严格的读写冲突检测实现完整的 NoConflict 公理，使事务的并发执行结果等价于某一串行执行序列，从而进一步消除写偏斜异常。

